\section*{Kurzfassung}
\label{sec:kurzfassung}
\addcontentsline{toc}{section}{\nameref{sec:kurzfassung}}
Der Kurs Algorithmen und Datenstrukturen 2 am Institut für Pervasive Computing hat traditionell statische Folien verwendet, um Konzepte wie Graphentheorie und Algorithmusvisualisierung zu vermitteln. Obwohl diese Methode zur Wissensvermittlung effektiv ist, bietet sie nur begrenzte Möglichkeiten für Studierende, interaktiv mit dem Material zu arbeiten, und für Lehrende, die Konzepte dynamisch zu demonstrieren.

Graphentheorie und Algorithmen sind grundlegende Themen der Informatik, und ihre dynamische Natur macht sie ideal für interaktive Erkundungen. Um dies zu ermöglichen, wird in dieser Arbeit eine webbasierte Anwendung vorgestellt, die sowohl Studierenden als auch Lehrenden erlaubt, eigene Graphen zu erstellen und Algorithmen intuitiv und praxisnah zu simulieren.

Diese Arbeit beschreibt den Entwurf, die Entwicklung und die Testphasen der Anwendung mit einem besonderen Fokus auf Benutzerfreundlichkeit und Funktionen, die speziell auf den Einsatz im Bildungsbereich zugeschnitten sind. Das Ergebnis ist ein praktisches Werkzeug, das das Lernen fördert, indem es Studierenden ermöglicht, Konzepte experimentell zu erforschen, und Lehrende dabei unterstützt, Algorithmen ansprechender und flexibler zu demonstrieren.